\subsection{Семинары «UML» (4 очных часа)}

Разработайте документ в текстовом процессоре, в котором полноценно 
описывается результат, полученный в предыдущей практической работе, с точки зрения
объектно-ориентированной архитектуры. Утверждения сопровождайте UML-диаграммами,
которые:
\begin{itemize}
    \item должны пояснять ваш текст;
    \item не должны быть перегруженными (в каждой диаграмме присутствуют только
    те классы и методы, которые поясняют описываемый аспект);
    \item должны строго соответствовать спецификации UML 2.5.1 (обращайте
    внимание на корректное использование стрелок).
\end{itemize}

В работе предполагается, что вы будете использовать:
\begin{itemize}
    \item диаграммы классов;
    \item диаграммы последовательности.
\end{itemize}